% consolidated.tex -- Merged paper: phys3 (sBootstrap) + phys4 (family symmetry)
% CRITICAL: mass = square of charge. Never "square root of mass."
\documentclass[12pt,a4paper]{article}

\usepackage{amsmath,amssymb,amsfonts,amsthm}
\usepackage{graphicx}
\usepackage{hyperref}
\usepackage{booktabs}
\usepackage{multirow}
\usepackage{geometry}
\usepackage{xcolor}
\geometry{margin=2.5cm}

\newtheorem{theorem}{Theorem}
\newtheorem{proposition}[theorem]{Proposition}
\newtheorem{corollary}[theorem]{Corollary}
\newtheorem{lemma}[theorem]{Lemma}

\newcommand{\SU}{\mathrm{SU}}
\newcommand{\UU}{\mathrm{U}}
\newcommand{\SO}{\mathrm{SO}}
\newcommand{\Sp}{\mathrm{Sp}}
\newcommand{\Nf}{N_{\!f}}
\newcommand{\Nc}{N_{\!c}}
\newcommand{\Tr}{\mathrm{Tr}}
\newcommand{\adj}{\mathrm{adj}}
\newcommand{\diag}{\mathrm{diag}}
\newcommand{\keV}{\,\mathrm{keV}}
\newcommand{\MeV}{\,\mathrm{MeV}}
\newcommand{\GeV}{\,\mathrm{GeV}}
\newcommand{\TeV}{\,\mathrm{TeV}}
\newcommand{\ie}{i.e.\@}
\newcommand{\eg}{e.g.\@}
\newcommand{\zch}{\mathfrak{z}}

\title{From mesinos and quarkinos to the Standard Model:\\
  the sBootstrap as a structural EFT from Seiberg confinement,\\
  charge--mass duality, and M-theory}

\author{A.\ Rivero\thanks{Corresponding author. \texttt{rivero@unizar.es}}}
\date{Draft --- \today}

\begin{document}
\maketitle

\begin{abstract}
We develop a unified framework in which Standard Model fermions
are identified with composite fermions of a confining
$\SU(3)_F$ family gauge theory: leptons are mesinos
(fermionic partners of the meson superfield $M^i{}_j = Q^i\tilde{Q}_j$),
and quarks are quarkinos arising through partial compositeness
in the complementary Higgs phase.
The confined spectrum naturally produces a doublet
$(M_i, \adj(M)_i)$ per family with the mathematical structure
of an $\mathcal{N}{=}2$ hypermultiplet, emergent from confinement
plus isospin.
Breaking this doublet structure to $\mathcal{N}{=}1$---by giving
a mass to the auxiliary adjoint---generates electroweak isospin:
the meson doublet $M^{1,2}$ becomes $H_d$ and the adjugate doublet
$\adj(M)^{1,2}$ becomes $H_u$, producing two Higgs doublets,
Yukawa couplings, and the up/down mass splitting from a single
composite field.

Having established the identification, we observe that the
mass spectrum of the extended meson sector of the Standard Model
exhibits a precise algebraic regularity.
The signed mass charge $\zch_k = z_0 + z_k$ (with
$m_k = \zch_k^2$, $\sum z_k = 0$, and $\langle z_k^2\rangle = z_0^2$)
satisfies the energy-balance condition $K = 2/3$ as a trace
identity of the $\UU(3)$ charge operator, independent of the
Cartan direction.
This regularity appears in the pseudoscalar mesons:
the seed triple $(0, \pi, D_s)$ to $0.18\%$ and the bloomed
triple $(-\pi, D_s, B)$ to $0.10\%$; in the charged leptons
$(e, \mu, \tau)$ to $0.001\%$; in the inverse down quarks
$K^{-1}(d,s,b)$ to $0.22\%$; and in the quark chain
$(t,b,c)$--$(b,c,-s)$--$(c,s,0)$--$(s,0,m_u{+}m_d)$.
The O'Raifeartaigh model produces the seed spectrum exactly
when the coupling ratio satisfies $gv/m = \sqrt{3}$.

The electroweak gauge bosons are accommodated through a
Casimir eigenvalue equation $x^2 + s(s{+}1)\,x - s(s{+}1) = 0$
that predicts $\sin^2\theta_W = 0.22310$ ($0.4\sigma$ from PDG)
and $M_W = 80.374\GeV$ ($0.39\sigma$).
A Diophantine system on the scalar composite content uniquely
selects $N = 3$ generations, with the composites filling
$\mathbf{24} \oplus \mathbf{15} \oplus \overline{\mathbf{15}}$
of $\SU(5)_f$---the representation content naturally contained
in the $\SO(32)$ adjoint but blocked by $E_8 \times E_8$.
In the M-theory lift, mesons are M2-branes ($C_3$ sector)
and baryons are M5-branes ($C_6$ sector), with the moduli-space
involution $M_i \leftrightarrow \Lambda^4/M_i$ realised as
electromagnetic duality of the $C$-field.
The membrane dimensionality $p = 2$ is the unique value giving
$K = 2/3$, connecting the charge-squared mass law to the
geometry of M-theory.
\end{abstract}

\tableofcontents
\newpage

%======================================================================
\section{Introduction: the sBootstrap philosophy}
\label{sec:intro}
%======================================================================

The Standard Model contains 13~free parameters in the flavour
sector---six quark masses, three charged-lepton masses, three CKM
angles, and a CP phase---with no known relations among them.
This paper develops a programme that addresses this problem not by
postulating new symmetries on the Yukawa matrices, but by
\emph{identifying the observed fermions with composite fermions}
of a confining gauge theory, in the spirit of nuclear
democracy~\cite{Chew:1962} and hadronic
supersymmetry~\cite{Miyazawa:1966,CattoGursey:1985}.

The foundational observation is the \emph{sBootstrap}
relation~\cite{Rivero:2024epjc}:
\begin{equation}
  m_\pi^2 - m_\mu^2 \;\approx\; f_\pi^2\,,
  \label{eq:pimu}
\end{equation}
where $f_\pi \approx 92.2\MeV$ is the pion decay constant.
Numerically, $m_\pi^2 - m_\mu^2 = 8\,316\MeV^2$ versus
$f_\pi^2 = 8\,501\MeV^2$ (a $2.2\%$ discrepancy).
This is the Volkov--Akulov mass formula for a
pseudo-Goldstino---the muon---whose scalar partner---the pion---is
a pseudo-Goldstone boson.
If taken seriously, Eq.~\eqref{eq:pimu} implies that
the pion and muon sit in the same $\mathcal{N}{=}1$ supermultiplet,
broken at scale $f_\pi$.

This identification has a precise realisation in
$\mathcal{N}{=}1$ $\SU(3)$ SQCD with $\Nf = \Nc = 3$ flavours.
The confined meson superfield $M^i{}_j = Q^i\tilde{Q}_j$
contains both a complex scalar (the meson) and a Weyl fermion
(the \emph{mesino}).
The sBootstrap postulate is:
\begin{quote}
\emph{Mesinos are leptons.  The nine fermionic components
of the $3 \times 3$ meson matrix furnish the SM lepton sector:
$3L_i + 3e^c_i$ (three doublets plus three singlets).}
\end{quote}
The scalar components of the same superfield provide the
Higgs doublets and sleptons.

The quark sector arises through a complementary mechanism.
In the Higgs phase of the same theory---smoothly connected to
the confining phase by the Fradkin--Shenker
theorem~\cite{FradkinShenker:1979}---the ``dressed quarks''
$\psi_M^{ij} \sim \tilde{v}_j \psi_Q^i + v_i \psi_{\tilde{Q}}^j$
are the same fermions as the mesinos.
Elementary quarks acquire mass through partial
compositeness~\cite{Kaplan:1991dc}: they mix linearly with
composite operators carrying the same quantum numbers.
The diquark partners needed to complete the SUSY multiplets
fill the symmetric $\mathbf{15} + \overline{\mathbf{15}}$ of
$\SU(5)_f$---a representation forbidden for point-like $qq$
composites by a kinematic theorem (the Pauli obstruction),
but naturally present in the $\SO(32)$ string adjoint.

The paper is organised to follow the logical flow of the
programme:
\begin{enumerate}
  \item \textbf{Sections~\ref{sec:composites}--\ref{sec:seiberg}}:
    The composite identification (quarks as quarkinos, leptons as
    mesinos) and the exact Seiberg dynamics that governs the
    confined spectrum.

  \item \textbf{Section~\ref{sec:N2}}: The emergent
    $\mathcal{N}{=}2$ doublet structure and its breaking to
    $\mathcal{N}{=}1$, which generates electroweak isospin, the
    two Higgs doublets, and the Yukawa structure.

  \item \textbf{Sections~\ref{sec:charges}--\ref{sec:meson_koide}}:
    The charge-squared mass law $m = \zch^2$, the energy-balance
    condition, and the observation that extended SM mesons
    satisfy the same algebraic regularity.

  \item \textbf{Section~\ref{sec:inverse}}: The Seiberg seesaw,
    inverse Koide, and the full chain of direct, inverse, and
    waterfall formulae.

  \item \textbf{Section~\ref{sec:WZ}}: The role of the $W$ and $Z$
    bosons: the Casimir eigenvalue equation and $\sin^2\theta_W$.

  \item \textbf{Sections~\ref{sec:SO32}--\ref{sec:Mtheory}}:
    Generation uniqueness, the $\SO(32)$ embedding, and the
    M-theory lift.
\end{enumerate}

Throughout, the mass of a particle is the \emph{square} of a
signed charge $\zch$: $m = \zch^2$.
The charge $\zch = z_0 + z_k$ can be positive or negative;
the mass is always non-negative.
This is not a notational convention but a physical statement:
the fundamental variable is the charge, and the mass is its
norm-square---the field energy stored in the gauge flux of a
confined composite state.


%======================================================================
\section{Quarks as quarkinos, leptons as mesinos}
\label{sec:composites}
%======================================================================

\subsection{The Miyazawa superalgebra}

The idea that mesons and baryons might be related by a
fermionic symmetry predates conventional supersymmetry.
Miyazawa~\cite{Miyazawa:1966} proposed in 1966 a graded Lie
algebra connecting mesons ($q\bar{q}$, bosons) to baryons
($qqq$, fermions), with the supercharge carrying colour
and flavour indices.
Catto and G\"ursey~\cite{CattoGursey:1985} developed this into
an algebraic framework where the diquark $[qq]$ serves as the
SUSY partner of the antiquark $\bar{q}$, so that the baryon
$q[qq]$ is the fermionic partner of the meson $q\bar{q}$.

In this picture, every quark has a diquark partner, and every
meson has a baryon partner.
The sBootstrap~\cite{Rivero:2024epjc} sharpens this: if the
SUSY-breaking scale is $f_\pi$ rather than the TeV scale, then
the meson--lepton pairing is \emph{literal}: the pion and muon
are superpartners, split by chiral symmetry breaking.

\subsection{The gauge-theoretic realisation}

The natural arena for this identification is $\mathcal{N}{=}1$
$\SU(3)_F$ SQCD with $\Nf = \Nc = 3$ flavours of chiral
superfields $Q^i$, $\tilde{Q}_j$ ($i,j = 1,2,3$).
Above the confinement scale $\Lambda_F$, the UV gauge group is
\begin{equation}
  G_{\rm UV} = \SU(3)_F \times \SU(3)_c \times \SU(2)_L
    \times \UU(1)_Y\,,
  \label{eq:GUV}
\end{equation}
where $\SU(3)_F$ is the family (preon-colour) gauge group that
confines at~$\Lambda_F$, while the SM gauge factors are weakly
gauged spectators.

The preons carry both family colour and electroweak quantum
numbers.
Embedding $\SU(2)_L \times \UU(1)_Y \subset \SU(3)_L$ via
the branching rule
$\mathbf{3}_L \to \mathbf{2}_{-1/2} \oplus \mathbf{1}_{+1}$,
the preon field content is:
\begin{center}
\begin{tabular}{lcccc}
\toprule
Field & $\SU(3)_F$ & $\SU(3)_c$ & $\SU(2)_L$ & $\UU(1)_Y$ \\
\midrule
$Q^{A}_a$ ($A=1,2$) & $\mathbf{3}$ & $\mathbf{1}$ & $\mathbf{2}$ & $-1/2$ \\
$Q^{3}_a$ & $\mathbf{3}$ & $\mathbf{1}$ & $\mathbf{1}$ & $+1$ \\
$\tilde{Q}^a{}_i$ & $\bar{\mathbf{3}}$ & $\mathbf{1}$ & $\mathbf{1}$ & $0$ \\
\bottomrule
\end{tabular}
\end{center}

After $\SU(3)_F$ confinement, the composite meson superfield
decomposes as:
\begin{equation}
  M^\alpha{}_i
  = \underbrace{M^{1,2}{}_i}_{(\mathbf{2}_{-1/2})}
  \;\oplus\;
  \underbrace{M^{3}{}_i}_{(\mathbf{1}_{+1})}\,.
  \label{eq:meson_decomp}
\end{equation}

The central sBootstrap identification is then:
\begin{equation}
  \boxed{
  \psi_{M^{1,2}{}_i} \equiv L_i\,,
  \qquad
  \psi_{M^{3}{}_i} \equiv e^c_i\,.
  }
  \label{eq:mesinos_leptons}
\end{equation}
The nine mesino fermions furnish the SM lepton sector:
three lepton doublets plus three charged singlets,
with anomaly cancellation guaranteed by 't~Hooft
matching~\cite{CsakiSchmaltzSkiba:1997}.

The scalar components of the same superfields provide:
\begin{center}
\begin{tabular}{llll}
\toprule
SQCD object & $\SU(2)_L \times \UU(1)_Y$ & sBootstrap role \\
\midrule
$M^{A}{}_i$ & $(\mathbf{2}, -1/2)$ & $L_i$ mesinos, $H_d$-type scalars \\
$M^{3}{}_i$ & $(\mathbf{1}, +1)$ & $e^c_i$ mesinos, charged scalars \\
$\adj(M)_A{}^{i}$ & $(\mathbf{2}, +1/2)$ & $H_u$-type composite operator \\
$B, \tilde{B}$ & $(\mathbf{1}, 0)$ & singlet baryonic superfields \\
\bottomrule
\end{tabular}
\end{center}

\subsection{The quark sector: diquarkinos and complementarity}

Quarks are \emph{not} mesinos---they carry colour.
In the sBootstrap, quarks are identified with effective coloured
multiplets that arise through Higgs--confinement
complementarity~\cite{FradkinShenker:1979}.
For $\SU(3)$ with $\Nf = 3$ fundamentals, the centre
$\mathbb{Z}_3$ is fully screened, and the Higgs and confinement
regimes are analytically connected: no local order parameter
distinguishes them.

In the Higgs phase, giving the squarks a diagonal VEV breaks
all 8~generators of $\SU(3)_F$.
After Higgsing, $18 - 8 = 10$ light chiral multiplets remain,
matching exactly the confined description:
$9~\text{mesons} + 2~\text{baryons} - 1~\text{constraint} = 10$.
The ``mesino'' in the confined description is the
\emph{same fermion} as the ``dressed quark'' in the Higgs phase.

The diquark partners $D \sim \wedge^2 Q$ needed to complete the
SUSY multiplets encounter a \textbf{bare diquarkino obstruction}:
requiring the meson doublet to furnish $(L_i, e^c_i)$ forces
the diquarkino singlet to carry hypercharge $Y = -4/3$ rather
than $+1/3$ or $-2/3$.
The resolution~\cite{Rivero:2024epjc} completes the flavour
sector into $\mathbf{15} \oplus \overline{\mathbf{15}}
\oplus \mathbf{24}$ of $\SU(5)_f$, with quarkino fields as
effective coloured multiplets of this completion.

A kinematic theorem shows that \emph{no} $J = 0$,
colour-$\bar{3}$ diquark can have symmetric flavour---regardless
of orbital angular momentum---so the sBootstrap diquark in the
$\mathbf{15}$ cannot be a $qq$ composite.
It must be a fundamentally string-theoretic state, as we discuss
in Section~\ref{sec:SO32}.


%======================================================================
\section{Seiberg confinement and the moduli space}
\label{sec:seiberg}
%======================================================================

\subsection{The quantum-modified constraint}

For $\Nf = \Nc = 3$, the theory confines with a
quantum-modified moduli space.
The infrared degrees of freedom---mesons $M^i{}_j$, baryon $B$,
antibaryon $\tilde{B}$---satisfy the exact quantum constraint
\begin{equation}
  \det M - B\tilde{B} = \Lambda^6\,,
  \label{eq:constraint}
\end{equation}
enforced by a Lagrange multiplier $X$ through the superpotential
\begin{equation}
  W = X\bigl(\det M - B\tilde{B} - \Lambda^6\bigr)
    + \sum_i m_i\,M^i{}_i\,.
  \label{eq:Wtot}
\end{equation}
The mass deformations $m_i$ break the flavour symmetry.
The $F$-term equations on the mesonic branch
($B = \tilde{B} = 0$) give the central result:
\begin{equation}
  \langle M_i \rangle = \frac{\Lambda^2 P^{1/3}}{m_i}\,,
  \qquad P = m_1 m_2 m_3\,.
  \label{eq:vacuum}
\end{equation}
\emph{Heavy preons produce light mesons, and vice versa.}
This is the Seiberg seesaw---the inverse relation that will
connect the charged-lepton and down-quark sectors.

\subsection{The adjugate mechanism}

The adjugate of the meson matrix,
$\adj(M)_i = \prod_{j \neq i} M_j$, is a
\emph{polynomial} (quadratic) in the meson entries---holomorphic,
no division required.
On the constraint surface $\det M = \Lambda^6$:
\begin{equation}
  \adj(M) = \Lambda^6\,M^{-1}\,,
  \label{eq:adjugate}
\end{equation}
so the adjugate provides a \emph{holomorphic realisation of the
matrix inverse}.
From the $F$-term equations:
\begin{equation}
  \adj(M)_i = -\Lambda^3\,m_i\,.
  \label{eq:Fterm}
\end{equation}
Thus the adjugate is proportional to the UV mass matrix,
while $M \propto 1/m$.
This gives two natural Yukawa operators at the confinement scale:
\begin{align}
  W_{\rm down} &\supset y_d\,q_L\,d^c\,\frac{M^{1,2}}{\Lambda}\,,
  & [H_d\text{-type}]\,,
  \label{eq:W_down} \\
  W_{\rm up} &\supset y_u\,q_L\,u^c\,
    \frac{\adj(M)^{1,2}}{\Lambda^3}\,,
  & [H_u\text{-type}]\,.
  \label{eq:W_up}
\end{align}
Species coupling through $M$ acquire masses
$m_{\rm phys} \propto 1/m_i$ (inverse hierarchy);
species coupling through $\adj(M)$ acquire
$m_{\rm phys} \propto m_i$ (direct hierarchy).
There are no independent Yukawa parameters: the mass hierarchy
\emph{is} the family charge hierarchy.


%======================================================================
\section{The emergent $\mathcal{N}{=}2$ doublet and electroweak isospin}
\label{sec:N2}
%======================================================================

\subsection{The composite doublet}

The confined spectrum contains two composite operators per
family: $M_i$ (meson) and $\adj(M)_i$ (adjugate).
Their mass eigenvalues have opposite charge dependence:
\begin{equation}
  \langle M_i \rangle \propto 1/m_i\,,
  \qquad
  \adj(M)_i \propto m_i\,.
  \label{eq:doublet}
\end{equation}
This pairing has the mathematical structure of an
$\mathcal{N}{=}2$ hypermultiplet: the two components are
rotated by an $\SU(2)_R$ that, in the $\mathcal{N}{=}2$ language,
exchanges the scalar and its dual.
Crucially, this doublet structure is \emph{emergent}---it arises
from confinement plus isospin, not from a UV $\mathcal{N}{=}2$
theory that is subsequently broken.

\subsection{From $\mathcal{N}{=}2$ to $\mathcal{N}{=}1$:
  the old trick}

The ``old trick'' of getting electroweak isospin from extended
SUSY works as follows.
Introduce an auxiliary adjoint chiral field $\Phi$ with
mass $\mu_\Phi$ and the standard $\mathcal{N}{=}2$ coupling
$W \supset \sqrt{2}\,g\,\tilde{Q}\,\Phi\,Q$.
Integrating out $\Phi$ via its $F$-term gives an effective
$\mathcal{N}{=}1$ superpotential with a quartic remnant:
\begin{equation}
  W_{\rm eff} = -\frac{g^2}{\mu_\Phi}
    \left[\Tr(Q\tilde{Q})^2
    - \frac{1}{3}\bigl(\Tr Q\tilde{Q}\bigr)^2\right]
    + m_i\,Q^i\tilde{Q}_i\,.
  \label{eq:Weff}
\end{equation}
As $\mu_\Phi \to \infty$ (decoupling), the quartic vanishes and
one recovers pure $\mathcal{N}{=}1$ SQCD.
At finite $\mu_\Phi$, the quartic couples the direct and inverse
branches.

The key structural consequence: the breaking
$\mathcal{N}{=}2 \to \mathcal{N}{=}1$ splits the hypermultiplet
$(M_i, \adj(M)_i)$ into two separate chiral superfields with
\emph{different} electroweak quantum numbers:
\begin{itemize}
  \item $M^{1,2}{}_i \sim (\mathbf{2}, -1/2)$:
    the down-type Higgs doublet $H_d$.
  \item $\adj(M)^{1,2}{}_i \sim (\mathbf{2}, +1/2)$:
    the up-type Higgs doublet $H_u$.
\end{itemize}
\textbf{Both MSSM Higgs doublets emerge from a single meson
matrix.}
The $\adj(M)$ is not an independent superfield but a holomorphic
polynomial in $M$, constrained by the quantum constraint
$\det M = \Lambda^6$.
At the superpotential level, this constrains $\tan\beta$:
$\langle\adj(M)_i\rangle / \langle M_i\rangle
= \Lambda^6 / \langle M_i\rangle^2$ is fixed by the quantum
constraint (subject to K\"ahler corrections).

The pairing of a direct mass ($\adj(M)_i \propto m_i$) with
an inverse mass ($M_i \propto 1/m_i$) within each family is
the composite manifestation of \textbf{electroweak isospin}:
up-type and down-type quarks in each generation acquire masses
from different faces of the same meson matrix.

\subsection{The CST alternative}

An alternative construction, due to Csaki, Shirman, and
Terning~\cite{Csaki:2011xn}, derives $\SU(2)_L$ dynamically
as the magnetic dual gauge group of a confining theory.
Consider an electric theory $\SU(N_e)$ with $\Nf$ flavours;
Seiberg duality gives a magnetic theory with gauge group
$\SU(\Nf - N_e)$.
Setting $\Nf - N_e = 2$, the magnetic gauge group is
$\SU(2)$, identified with $\SU(2)_L$.
The magnetic superpotential
$W_{\rm mag} = q\,\mathcal{M}\,\tilde{q}/\mu$ generates
Yukawa couplings between magnetic quarks and the meson Higgs.
This costs two extra parameters ($\Lambda_{\rm CST}$ and the
electric gauge coupling) that do not enter any numerical
prediction; the branching-rule embedding of the main text
achieves the same Higgs identification with zero additional
parameters.


%======================================================================
\section{The charge-squared mass law}
\label{sec:charges}
%======================================================================

\subsection{Charges from $\UU(3)$ normalisation}

Let the three generations transform as the fundamental
$\mathbf{3}$ of $\UU(3)_F = \SU(3)_F \times \UU(1)$.
The Cartan charges satisfy
\begin{equation}
  \sum_{i=1}^3 z_i = 0\,,
  \qquad
  \sum_{i=1}^3 z_i^2 = \frac{1}{2}\,,
  \label{eq:su3}
\end{equation}
with the $\UU(1)$ charge fixed by canonical normalisation:
$z_0 = 1/\sqrt{6}$.
The total signed charge of the $i$-th species is
\begin{equation}
  \zch_i \;\equiv\; z_0 + z_i\,,
  \label{eq:charge}
\end{equation}
and the physical mass is
\begin{equation}
  \boxed{m_i \;=\; \mu\,\zch_i^2\,,}
  \label{eq:massansatz}
\end{equation}
with $\mu$ a generation-independent scale.
The charge can be positive or negative; the mass is always
non-negative.

\subsection{Why the square: four routes}

The quadratic dependence $m \propto \zch^2$ has multiple
independent motivations:

\paragraph{(i) Field energy.}
By Gauss's law, the chromoelectric field of a static preon
with family charge $\zch_i$ is $|\mathbf{E}_i| \propto \zch_i$.
Since $\mathcal{L} \supset -(1/4g^2)F_{\mu\nu}^2$, the stored
field energy is
$U_i = (2g^2)^{-1}\int d^3x\,\mathbf{E}_i^2 \propto \zch_i^2$---quadratic in charge, like a capacitor $E = Q^2/(2C)$.

\paragraph{(ii) Bilinear composites.}
The meson $M^i{}_i = Q^i\tilde{Q}_i$ is a bilinear in the preon
fields.
If each preon contributes a factor proportional to $\zch_i$,
the composite VEV scales as $\zch_i^2$.
This is the mechanism identified by Koide~\cite{Koide:1990}:
any mass arising through \emph{two} insertions of a VEV
(a seesaw or radiative loop) gives $m \propto v^2$, which is
the charge-squared structure.

\paragraph{(iii) Membrane area.}
In the M-theory lift, the meson is an M2-brane stretched
between D6-branes.
If both worldvolume directions scale linearly with $\zch_i$,
the BPS mass is $m \propto T_2 \times \text{Area}
\propto \zch_i^2$.
This is the unique brane dimensionality that produces
$K = 2/3$ (Section~\ref{sec:Mtheory}).

\paragraph{(iv) Uniqueness theorem.}
For monomial mass laws $m_i = c\,(z_0 + z_i)^n$, the
generalised ratio
$K_n = \sum (z_0{+}z_i)^n / (\sum (z_0{+}z_i)^{n/2})^2$
is independent of the Cartan direction $\alpha$ if and only
if $n = 2$ (besides the trivial $n = 0$).
For $n \geq 3$, the cubic Casimir
$\sum z_i^3 \propto \cos 3\alpha$ introduces
direction-dependent terms.

\subsection{The energy-balance condition}

For three charges $\zch_k = z_0 + z_k$ with $\sum z_k = 0$:
\begin{equation}
  K \;=\; \frac{\sum \zch_k^2}{(\sum \zch_k)^2}
  \;=\; \frac{3z_0^2 + \sum z_k^2}{9z_0^2}
  \;=\; \frac{1}{3} + \frac{\langle z_k^2\rangle}{3z_0^2}\,.
  \label{eq:K_derivation}
\end{equation}
Therefore:
\begin{equation}
  \boxed{K = \frac{2}{3}
  \quad\Longleftrightarrow\quad
  \langle z_k^2\rangle = z_0^2\,.}
  \label{eq:energy_balance}
\end{equation}
The energy balance condition says: the mean energy of the
perturbations equals the energy of the unperturbed charge.
This is neither perturbative ($\sigma \ll \bar{\zch}$,
$K \to 1/3$, all masses equal) nor dominant
($\sigma \gg \bar{\zch}$, $K \to 1$, one mass dominates).
The value $K = 2/3$ is the equipartition point.

One mass can vanish naturally: if
$\langle z_k^2 \rangle = z_0^2$, a perturbation $z_k$ can reach
$|z_k| \sim z_0$, making $\zch_k = z_0 + z_k \approx 0$.
A vanishing mass is the natural extreme of energy balance,
not fine-tuning.

\subsection{The shift parameter}

Promoting the fixed offset $z_0 = 1/\sqrt{6}$ to a free shift
parameter $s$, the Koide ratio on the direct branch becomes
\begin{equation}
  K(s) = \frac{1}{3} + \frac{1}{18 s^2}\,,
  \label{eq:K_of_s}
\end{equation}
independent of the Cartan angle $\alpha$.
The condition $K = 2/3$ selects $s = 1/\sqrt{6}$ uniquely.
Systematic fitting of all known tuples reveals that $s$ is
approximately \emph{universal}:
$s \in [0.33, 0.49]$ with a mean of $0.40 \approx z_0$,
while the Cartan angle $\alpha$ spans $4^\circ$--$37^\circ$
across sectors.
The universality of $s$---not $\alpha$---is the structural
content of the regularity.


%======================================================================
\section{The extended meson spectrum}
\label{sec:meson_koide}
%======================================================================

\subsection{Seed triples}

For a triple $(0, \zch_a, \zch_b)$ with one massless member,
the energy balance condition requires:
\begin{equation}
  \frac{\zch_a}{\zch_b} = 2 - \sqrt{3} \approx 0.2680\,,
  \label{eq:seed_ratio}
\end{equation}
a \emph{universal} prediction: any charge triple with one zero
member has the same charge ratio between the two massive members,
regardless of their absolute scale.

\paragraph{The quark seed: $(0, s, c)$.}
The strange and charm quark charges give
$\zch_s / \zch_c = 9.67 / 35.64 = 0.271$,
deviating from the seed ratio by $1.2\%$.
The Koide diagnostic gives $K = 0.664$ ($0.35\%$ from $2/3$).

\paragraph{The meson seed: $(0, \pi, D_s)$.}
The pion and $D_s$ meson charges give
$\zch_\pi / \zch_{D_s} = 11.81 / 44.37 = 0.266$,
deviating by $0.6\%$.
The Koide diagnostic gives $K = 0.668$ ($0.18\%$ from $2/3$).
This is the \textbf{meson counterpart} of the quark seed---both
involve the strange-charm sector.
The two seeds are nearly parallel in charge space ($0.3^\circ$):
the meson seed is the composite image of the quark seed.

\paragraph{The lepton triple as ``almost-seed.''}
The charged lepton triple $(e, \mu, \tau)$ has
$\zch_e = 0.715\MeV^{1/2}$, which is small but nonzero.
The leptons are not an exact seed---the electron, while very
light, is not massless.
But the triple is the same family: the near-critical charge
puts one member close to zero.

\subsection{The meson triple $(-\pi, D_s, B)$}

With the pion charge \emph{negative}, the triple satisfies:
\begin{equation}
  K(-\pi, D_s, B) = \frac{m_\pi + m_{D_s} + m_B}
  {(-\zch_\pi + \zch_{D_s} + \zch_B)^2}
  = 0.6674\,,
  \label{eq:meson_koide}
\end{equation}
deviating from $2/3$ by $0.10\%$---the second most precise
charge triple known (after the leptons).
The pion's negative charge has physical meaning: as a
pseudo-Goldstone boson, the pion approaches $\zch = 0$ from below
as $m_q \to 0$.
The sign encodes Goldstone nature.

Among $\binom{11}{3} \times 4 = 660$ combinations of the
eleven pseudoscalar mesons with sign patterns, only
$(-\pi, D_s, B)$ and the closely related $(-\pi, D_s, B_s)$
lie within $0.2\%$ of $K = 2/3$.
All six closest hits share the sign assignment
$(-\pi, D_{(s)}, B_{(s,c)})$: a negative pion paired with
charm- and bottom-containing mesons---the combination predicted
by the $\SU(5)_f$ flavour structure.

\subsection{The O'Raifeartaigh mechanism}

The seed-to-bloom transition has a precise SUSY-breaking
realisation.
The O'Raifeartaigh model with three chiral superfields and
superpotential $W = f\Phi_0 + m\Phi_1\Phi_2 + g\Phi_0\Phi_1^2$
produces a fermionic mass spectrum $(0, m_-, m_+)$ with
$m_\pm = (\sqrt{g^2v^2 + m^2} \pm gv)$.
The charge ratio $\zch_- / \zch_+ = 2 - \sqrt{3}$ (the seed
condition) requires:
\begin{equation}
  \boxed{\frac{gv}{m} = \sqrt{3}\,,}
  \label{eq:OR_seed}
\end{equation}
a \emph{single algebraic constraint} on the superpotential
parameters.
The Koide condition on the seed is therefore one equation
relating the cubic coupling, the pseudo-modulus VEV, and the
bilinear mass parameter.

The bloom from seed to full triple is a phase rotation
$\delta: 3\pi/4 \to \delta_{\rm phys}$ in the Koide
parametrisation, at fixed scale---mass redistributes among
the three eigenvalues without changing their sum.
The O'Raifeartaigh seed is an \emph{unstable fixed point} of
the R-breaking perturbation: any nonzero R-breaking pushes
$K$ away from $2/3$.
The bloom must therefore be a nonperturbative rearrangement.


%======================================================================
\section{The Seiberg seesaw and the inverse sector}
\label{sec:inverse}
%======================================================================

\subsection{The dual Koide on $(d, s, b)$}

At the Seiberg vacuum, $\langle M^j{}_j \rangle \propto 1/m_j$:
heavy quarks produce light mesons and vice versa.
This prompts the question: does the inverse-mass spectrum
$\{1/m_i\}$ satisfy the energy-balance condition?
Evaluating for the pure down-type quarks:
\begin{equation}
  K(1/m_d,\, 1/m_s,\, 1/m_b) = 0.6652\,,
  \label{eq:dual_koide}
\end{equation}
a $0.22\%$ deviation from $2/3$---better than any direct quark
triple.
No other quark triple comes close under the Seiberg seesaw
(the next best, $(d,s,c)$, deviates by $4.2\%$).

This ``dual Koide'' connects the UV mass hierarchy to the IR
meson condensates: the Seiberg constraint maps the down-type
quark masses into a near-energy-balance spectrum of meson VEVs.
The dual charges $\xi_j = 1/\zch_j$ are the charges of the
Seiberg meson VEVs---the magnetic Koide is an energy balance
on the meson condensate, not on quark masses.

A simple algebraic identity connects the seed to the seesaw:
$K(a, b) = K(1/a, 1/b)$ identically for any two masses.
The seed $(0, m_s, m_c)$ is thus the unique self-dual fixed point
of the seesaw: the electric and magnetic energy-balance conditions
coincide only when one mass vanishes.

\subsection{The interleaving theorem and the waterfall}

The two branches---direct ($a_i = \mu\,\zch_i^2$) and
inverse ($d_i = \nu/\zch_i^2$)---have opposite orderings:
$a_1 < a_2 < a_3$ while $d_1 > d_2 > d_3$.
This is the algebraic origin of the quark mass waterfall.

\begin{theorem}[Interleaving]
The six unmixed masses form the alternating ladder
$d_1 > a_3 > d_2 > a_2 > d_3 > a_1$
if and only if
$\max(\zch_1^2\zch_3^2, \zch_2^4) < \nu/\mu < \zch_2^2\zch_3^2$.
\end{theorem}

The SM quark masses satisfy this window condition.
The resulting 6-rung ladder alternates strictly between inverse
and direct entries, producing consecutive windows that mix both
branches.
The sign flip in $K(b,c,-s)$ signals the branch mismatch in the
mixed window.

The geometric mean $\sqrt{a_i \cdot d_i} = \sqrt{\mu\nu}$ is
family-independent---a direct consequence of $\det M = \Lambda^6$.

\subsection{The chain structure}

The quark predictions arise from a chain of overlapping triples:
\begin{equation}
  (t,b,c) - (b,c,s)^{\pm} - (c,s,0) - (s,0,m_u{+}m_d)\,,
  \label{eq:chain}
\end{equation}
split into two regimes.
On the \textbf{Fermi side} (above the SUSY-breaking scale):
$K(m_t, m_b, m_c) \approx 2/3$ (standard, $0.39\%$).
On the \textbf{chiral side} (below the breaking scale):
$m_c/m_s = (2{+}\sqrt{3})^2 \approx 13.93$ and
$m_s/(m_u{+}m_d) = (2{+}\sqrt{3})^2$.
The bridge tuple $(b,c,s)$ with a signed charge flip connects
the two sides.

The $v_0$-doubling relation
$\zch_b = 3\zch_s + \zch_c$ matches
PDG to $0.07\%$ ($m_b = 4177$ vs.\ $4180\MeV$).
This arises from a bion K\"ahler correction
$K_{\rm bion} \propto |\sum s_k \zch_k|^2$ from
monopole-instantons on $\mathbb{R}^3 \times S^1$.

\subsection{CKM predictions}

The self-dual charge spectrum has $R = (5+\sqrt{21})/2$,
the fundamental unit of $\mathbb{Z}[(1+\sqrt{21})/2]$.
The CKM elements:
\begin{center}
\begin{tabular}{lllll}
\toprule
Observable & Formula & Predicted & PDG 2024 & Status \\
\midrule
$|V_{us}|$ & $\sqrt{m_d/m_s}$ & $0.2236$ & $0.2243 \pm 0.0008$ & $0.9\sigma$ \\
$|V_{cb}|$ & $1/(5R)$ & $0.0417$ & $0.0422 \pm 0.0008$ & $0.6\sigma$ \\
$|V_{ub}|$ & $2\sqrt{2}/(7R^3)$ & $0.00367$ & $0.00368 \pm 0.00011$ & $0.1\sigma$ \\
$\delta_{\rm CP}$ & $\arctan\sqrt{R}$ & $65.4^\circ$ & $65.7 \pm 3.4^\circ$ & $0.1\sigma$ \\
$J$ & closed form & $3.04 \times 10^{-5}$ & $3.08 \times 10^{-5}$ & $0.3\sigma$ \\
\bottomrule
\end{tabular}
\end{center}
The CP phase is equivalent to
$\cos^2(2\delta) = N/(N{+}4) = N/(2N{+}1)$, which holds
only for $N = 3$---the same condition as discriminant matching.
The integers 5 and 7 are empirical pattern identifications;
derivations from the operator basis remain open.


%======================================================================
\section{The $W$ and $Z$ bosons}
\label{sec:WZ}
%======================================================================

\subsection{The Casimir eigenvalue equation}

Consider the eigenvalue equation
\begin{equation}
  x^2 + C_2\,x - C_2 = 0\,,
  \qquad C_2 = s(s{+}1)\,,
  \label{eq:casimir}
\end{equation}
where $C_2$ is the quadratic Casimir of $\SU(2)$ at spin $s$.
This equation is pinned by three structural constraints:
(i) no spin-independent mass term ($f_0 = 0$);
(ii) a massless eigenstate at $s = 0$ ($g_0 = 0$);
(iii) coefficient antisymmetry ($g_1 = -f_1$).

Evaluating at $s = 1/2$ ($C_2 = 3/4$) and $s = 1$ ($C_2 = 2$):
\begin{align}
  x_+(1/2) &= \tfrac{1}{8}(-3 + \sqrt{57})\,,
  \quad x_+(1/2)\,m^2 = M_W^2\,,\\
  x_+(1) &= -1 + \sqrt{3}\,,
  \quad x_+(1)\,m^2 = M_Z^2\,.
\end{align}
The electroweak ratio is
\begin{equation}
  R \;\equiv\; 1 - \frac{x_+(1/2)}{x_+(1)}
  = \frac{(\sqrt{19} - 3)(\sqrt{19} - \sqrt{3})}{16}
  \approx 0.22310\,,
  \label{eq:sin2}
\end{equation}
to be compared with
$\sin^2\theta_W^{\rm os} = 0.22320 \pm 0.00026$ ($0.4\sigma$).
The predicted $W$ mass is
$M_W^{\rm pred} = 80.374 \pm 0.002\GeV$ ($0.39\sigma$ from PDG).

\subsection{The fourth eigenvalue}

At $s = 1/2$, the second root gives
$|M(1/2, -)| = 122.4\GeV$, which is $2.3\%$ below
$M_H = 125.25\GeV$.
This eigenstate carries spin $1/2$ while the Higgs is spin $0$;
the near-degeneracy is noted as a numerical curiosity.

\subsection{The scale parameter}

Setting $x_+(1)\,m^2 = M_Z^2$ gives
$m = M_Z/\sqrt{\sqrt{3} - 1} = 106.6\GeV \approx m_\mu \times 10^3$.

\subsection{Origin and scheme dependence}

This observation was first made by H.~de~Vries (2004) and
subsequently studied in~\cite{Rivero:2005b}.
The comparison uses the on-shell definition
$\sin^2\theta_W = 1 - M_W^2/M_Z^2 = 0.22320$,
\textbf{not} the $\overline{\rm MS}$ value $0.23122$.
Until a dynamical mechanism selecting the on-shell scheme is
identified, the agreement should be regarded as suggestive but
scheme-dependent.


%======================================================================
\section{Generation uniqueness and the $\SO(32)$ embedding}
\label{sec:SO32}
%======================================================================

\subsection{The Diophantine system}

The sBootstrap postulates that every scalar of the
Supersymmetric Standard Model is a composite of two quarks
(diquark) or of a quark--antiquark pair (meson).
Matching composites to fermions, charge sector by charge sector:
\begin{align}
  rs &= 2N\,, \label{eq:rs}\\
  \frac{r(r{+}1)}{2} &= 2N\,, \label{eq:rr}\\
  r^2 + s^2 - 1 &= 4N\,. \label{eq:neutral}
\end{align}

\begin{theorem}[Generation uniqueness]
The system~\eqref{eq:rs}--\eqref{eq:neutral} has a unique
positive-integer solution:
$(N, r, s) = (3, 3, 2)$.
\end{theorem}

\noindent With $r + s = 5$ light quarks forming the fundamental
$\mathbf{5}$ of $\SU(5)_f$, the composites organise into
$\mathbf{24} \oplus \mathbf{15} \oplus \overline{\mathbf{15}}$.
The counting equation~\eqref{eq:rr} selects the
\emph{symmetric} product $\mathbf{15}$, not the antisymmetric
$\overline{\mathbf{10}}$---the Pauli obstruction for
point-like diquarks.

\subsection{Discriminant matching}

Two numbers arise naturally for $N$ generations:
the algebraic discriminant $D_{\rm alg} = (N{+}2)^2 - 4$
and the dynamical discriminant $D_{\rm dyn} = N(2N{+}1)$.
The condition $D_{\rm alg} = D_{\rm dyn}$ reduces to
$N^2 - 3N = 0$:
\begin{equation}
  \boxed{N = 3\,.}
  \label{eq:N3}
\end{equation}
At $N = 3$: $D = 21 = \dim \Sp(6) = \dim \adj\,\SO(7)$.

\subsection{The $\SO(32)$ adjoint}

The $\SO(32)$ gauge group of Type~I / heterotic string theory
admits the chain
$\SO(32) \supset \SU(16) \times \UU(1)
\supset \SU(5) \times \SU(3) \times \UU(1)^2$.
The adjoint $\mathbf{496} = \wedge^2(\mathbf{32})$ decomposes
to contain, in the neutral sector ($Q_A = Q_B = 0$):
$\mathbf{24} \oplus \mathbf{1}$, and in the matter sector
($|Q_A| = 2$):
$\mathbf{15} \oplus \overline{\mathbf{15}}
\oplus \mathbf{10} \oplus \overline{\mathbf{10}}
\oplus \mathbf{5} \oplus \bar{\mathbf{5}}$.

The sBootstrap representations $\mathbf{24}$ and
$\mathbf{15} + \overline{\mathbf{15}}$ appear naturally.
The remaining 442 states must be projected out by an
orbifold or GSO-type condition---an open problem.

\paragraph{Why $\SO(32)$ and not $E_8 \times E_8$.}
Every decomposition of $E_8$ under $\SU(5)$ produces the
antisymmetric $\mathbf{10}$, never the symmetric $\mathbf{15}$.
The obstruction is irreducible: $E_8$'s simply-laced root lattice
admits only antisymmetric tensor representations of $\SU(5)$.
By contrast, $\SO(32)$'s adjoint naturally contains the symmetric
$\mathbf{15}$ through
$\wedge^2(\mathbf{5} \otimes \mathbf{3})
\supset S^2(\mathbf{5}) \otimes \wedge^2(\mathbf{3})
= (\mathbf{15}, \bar{\mathbf{3}})$.
The sBootstrap \textbf{constrains the string embedding to the
Type~I / $\SO(32)$ corner} of the landscape.


%======================================================================
\section{M-theory lift}
\label{sec:Mtheory}
%======================================================================

\subsection{Mesons as M2-branes, baryons as M5-branes}

The Hanany--Witten brane construction~\cite{HananyWitten:1997}
realises $\SU(3)$ SQCD in Type~IIA using D4-branes between
NS5-branes with D6-branes providing flavours.
Lifted to M-theory, the configuration merges into a single
M5-brane wrapping a holomorphic curve $\Sigma$.

\begin{center}
\begin{tabular}{llll}
\toprule
Field theory & Type~IIA & M-theory & Mass type \\
\midrule
Meson $M^i{}_j = Q^i\tilde{Q}_j$ &
  open string (D6--D6) &
  \textbf{M2-brane} & Dirac \\
Baryon $B = \epsilon\,Q^3$ &
  wrapped D4 (vertex) &
  \textbf{M5-brane} & Majorana \\
\bottomrule
\end{tabular}
\end{center}

The moduli-space involution
$M_i \to \Lambda^4/M_i$ exchanges the $C_3$ and $C_6$
descriptions---the \emph{electromagnetic duality} of M-theory.

\subsection{Why $p = 2$: the membrane signature}

With $\UU(3)$ family symmetry, D6-brane positions are
$v_i = c\,\zch_i$.
An M2-brane stretched to the $i$-th D6-brane spans two
spatial directions, so its BPS mass is
$m_i \propto \text{Area}_i \propto \zch_i^2$.

For a general $p$-brane:
\begin{center}
\begin{tabular}{cccc}
\toprule
$p$ & Object & $m_i \propto$ & $K$ \\
\midrule
1 & string & $\zch_i$ & 0.478 \\
\textbf{2} & \textbf{M2-brane} & $\zch_i^2$ & \textbf{2/3} \\
3 & 3-brane & $\zch_i^3$ & 0.805 \\
5 & M5-brane & $\zch_i^5$ & 0.945 \\
\bottomrule
\end{tabular}
\end{center}
\textbf{Only $p = 2$ yields $K = 2/3$}: the energy-balance
condition is a signature of the membrane.

\subsection{The number 84}

In 11d SUGRA, the $C_3$ field has
$\binom{9}{3} = 84$ polarisations.
Under $\SU(9) \subset E_8$:
$248 = 80 + 84 + \overline{84}$.
In the SM with right-handed neutrinos, the 96~fermionic
degrees of freedom split as $84 + 12$ under dual
mass-protection symmetries.
A fourth uniqueness criterion:
$\binom{3N}{N} = (N{+}1) \cdot N(2N{+}1)$ holds only at $N = 3$,
giving $84 = 4 \times 21$.

\subsection{Strong CP and the involution as CP}

Identifying the moduli-space involution with CP gives a
Nelson--Barr structure: the superpotential is real in the
CP-symmetric basis, so $\arg\det m_q = 0$ at tree level.
The CKM phase arises from spontaneous CP violation.
Radiative corrections:
$\delta\bar\theta \sim (\alpha/4\pi)^2 J \sim 10^{-11}$,
at or below the experimental bound $|\bar\theta| < 10^{-10}$.


%======================================================================
\section{The explicit Lagrangian}
\label{sec:lagrangian}
%======================================================================

Below the family confinement scale $\Lambda_F$, the effective
theory is specified by:

\paragraph{K\"ahler potential.}
\begin{equation}
  K = \underbrace{\Tr(M^\dagger M) + |X|^2 + |B|^2
    + |\bar{B}|^2}_{K_{\rm can}}
  \underbrace{{} - \frac{|X|^4}{12\mu^2}}_{K_{\rm stab}}
  + \underbrace{\frac{\zeta^2}{\Lambda^2}
    \Bigl|\sum_{k} s_k \zch_k\Bigr|^2}_{K_{\rm bion}}\,,
  \label{eq:Kahler}
\end{equation}
where $K_{\rm stab}$ pins the pseudo-modulus at
$\langle X \rangle = \sqrt{3}\,\mu$ (the K\"ahler pole),
and $K_{\rm bion}$ enforces the $v_0$-doubling condition
$\zch_b = 3\zch_s + \zch_c$.

\paragraph{Superpotential.}
\begin{equation}
  W = \underbrace{\Tr(\hat{m}\,M)}_{\rm chiral~mass}
    + \underbrace{X(\det M - B\bar{B} - \Lambda^6)}_{\rm Seiberg}
    + \underbrace{c_3\frac{(\det M)^3}{\Lambda^{18}}}_{\rm 3-inst.}
    + \underbrace{y_c H_u M^c{}_c + y_b H_d M^b{}_b
      + y_t H_u Q^t\bar{Q}_t}_{\rm Yukawa}
    + \underbrace{\lambda_S S H_u \cdot H_d
      + \frac{\kappa}{3}S^3}_{\rm NMSSM}\,.
  \label{eq:W}
\end{equation}
The three-instanton term $(\det M)^3/\Lambda^{18}$ generates a
$\cos(9\delta)$ potential whose minima include the observed lepton
Koide phase to $0.9\sigma$ in $m_\tau$.
The NMSSM singlet $S$ is a separate field from the Lagrange
multiplier $X$ (dimensional analysis:
$\lambda_{\rm Higgs}/\lambda_{\rm max} \sim 10^5$).
At $\tan\beta = 1$, the Higgs mass is
$m_h = \lambda_S v/\sqrt{2} = 125\GeV$ for $\lambda_S = 0.72$.

\paragraph{Mixed triples require $\tan\beta = 1$.}
The quark Koide triples obligatorily mix up-type and down-type
quarks.
The translation from masses to Yukawa couplings depends on
$\tan\beta$: $K$ on masses equals $K$ on Yukawas only when
$\sin\beta = \cos\beta$, i.e., $\tan\beta = 1$.

\paragraph{Soft breaking.}
The soft terms arise from a spurion $\langle F_S \rangle = f_\pi^2$:
$V_{\rm soft} = f_\pi^2\,\Tr(M^\dagger M)$.
The supertrace $\mathrm{STr}[M^2] = 18\,f_\pi^2$ receives
contributions only from the nine soft meson masses.


%======================================================================
\section{SUSY breaking: $F$-terms are lepton masses}
\label{sec:susy_breaking}
%======================================================================

The SUSY-preserving vacuum is diagonal and does not break
$\SU(2)_L$.
When electroweak symmetry breaks (all three Higgs doublet
components acquire VEVs), the $F$-term system becomes
overconstrained: 18~real equations for 12~real variables.
SUSY \emph{necessarily} breaks, with
\begin{equation}
  |F_{2j}| = m_j\,, \qquad j = e, \mu, \tau\,.
  \label{eq:F_lepton}
\end{equation}
\textbf{The SUSY-breaking auxiliary fields \emph{are} the
lepton masses.}
The goldstino is predominantly the $\tau$-mesino, and
$m_{3/2} = |F|/(\sqrt{3}\,M_{\rm Pl})
\approx 4 \times 10^{-16}\GeV$.

The three Higgs doublets $H_i = M^{1,2}{}_i$ acquire
hierarchical VEVs $v_i \propto m_{\ell,i}$: the SM Higgs is
predominantly the $\tau$-generation doublet
($v_\tau/\sqrt{\sum v_i^2} = 0.998$).


%======================================================================
\section{Numerical results and status}
\label{sec:numerics}
%======================================================================

\subsection{Parameter counting}

The framework has four continuous inputs:
$k$ (overall scale, from $m_e$),
$\alpha$ (sector angle, from $m_\mu$),
$m_d$ and $m_s$ ($\overline{\rm MS}$ at $2\GeV$, from PDG).
From these, it produces:

\begin{table}[ht]
\centering\small
\begin{tabular}{lllll}
\hline
Observable & Method & Predicted & PDG 2024 & Status \\
\hline
$m_\tau$ & $K = 2/3$ & $1776.97\MeV$ & $1776.86 \pm 0.12$ & $0.9\sigma$ \\
$K(e,\mu,\tau)$ & trace identity & $2/3$ & $0.666661$ & $0.001\%$ \\
$K^{-1}(d,s,b)$ & Seiberg seesaw & $\approx 2/3$ & $0.6652$ & $0.22\%$ \\
$K(-\pi,D_s,B)$ & meson triple & $\approx 2/3$ & $0.6674$ & $0.10\%$ \\
$m_c$ & $(2{+}\sqrt{3})^2 m_s$ & $1301\MeV$ & $1273$ & $2.4\%$ \\
$m_u + m_d$ & zero-charge & $6.7\MeV$ & $6.8$ & $1.4\%$ \\
$|V_{us}|$ & $\sqrt{m_d/m_s}$ & $0.2236$ & $0.2243$ & $0.9\sigma$ \\
$|V_{cb}|$ & $1/(5R)$ & $0.0417$ & $0.0422$ & $0.6\sigma$ \\
$|V_{ub}|$ & $2\sqrt{2}/(7R^3)$ & $0.00367$ & $0.00368$ & $0.1\sigma$ \\
$\delta_{\rm CP}$ & $\arctan\sqrt{R}$ & $65.4^\circ$ & $65.7^\circ$ & $0.1\sigma$ \\
$\sin^2\theta_W$ & Casimir eq. & $0.22310$ & $0.22320$ & $0.4\sigma$ \\
$M_W$ & from above & $80.374\GeV$ & $80.369$ & $0.39\sigma$ \\
$m_b$ (bion) & $v_0$-doubling & $4177\MeV$ & $4183$ & $0.06\sigma$ \\
$m_t$ (chain) & $(t,b,c)$ & $\sim 171\GeV$ & $172.8$ & $1\%$ \\
\hline
\end{tabular}
\caption{Quantitative results.  $R = (5+\sqrt{21})/2$.}
\label{tab:results}
\end{table}

\subsection{Epistemological classification}

\begin{table}[ht]
\centering\small
\begin{tabular}{lll}
\hline
Claim & Status & Comment \\
\hline
$K = 2/3$ from $\UU(3)$ charges & Derived & trace identity \\
$n = 2$ uniqueness & Derived & binomial theorem \\
Seiberg transparency $(W_S)_{kk} = 0$ & Derived & det multilinearity \\
Involution $M \leftrightarrow \Lambda^4/M$ & Derived & moduli space \\
Interleaving theorem & Derived & opposite orderings \\
$N = 3$ uniqueness (Diophantine) & Derived & $(r-3)(r+1) = 0$ \\
$N = 3$ uniqueness (discriminant) & Conditional & equal quartic couplings \\
O'Raifeartaigh seed at $gv/m = \sqrt{3}$ & Derived & fermion eigenvalues \\
Pauli theorem (no symmetric $J{=}0$ diquark) & Derived & angular momentum \\
$\tan\beta = 1$ from mixed triples & Derived & $\sin\beta = \cos\beta$ \\
$\SO(32)$ but not $E_8 \times E_8$ & Derived & $\mathbf{15}$ vs.\ $\mathbf{10}$ \\
\hline
Lepton Koide $K = 2/3$ to $0.001\%$ & Observed & $2.8\sigma$ after LEE \\
Meson triple $(-\pi, D_s, B)$ to $0.10\%$ & Observed & $1.1\sigma$ after LEE \\
$v_0$-doubling: $m_b$ to $0.07\%$ & Observed & bion K\"ahler \\
Casimir $R = 0.22310$ vs.\ $\sin^2\theta_W$ & Observed & scheme-dependent \\
Dual Koide $K^{-1}(d,s,b) = 0.665$ & Observed & $0.22\%$ from $2/3$ \\
$|V_{cb}| = 1/(5R)$, integers 5, 7 & Empirical & no derivation yet \\
$\alpha_{\rm lep}$ origin & Open & K\"ahler modulus \\
\hline
\end{tabular}
\caption{Logical status of claims.}
\label{tab:status}
\end{table}


%======================================================================
\section{TeV-scale convergence}
\label{sec:TeV}
%======================================================================

Three TeV-range scales emerge with no TeV-scale input:
\begin{enumerate}
  \item The emergent adjoint mass
    $\mu_{\rm crit} \approx 4.5\TeV$ (below which the Koide
    vacuum ceases to exist).
  \item The composite stop mass $M_S(X_t = 0) \approx 4.6\TeV$
    required for $m_h = 125\GeV$ via the Coleman--Weinberg formula.
  \item The inverse Koide crossing scale
    $\mu_K \approx 1\TeV$ where $K^{-1}(d,s,b)$ passes through
    $2/3$ under 5-loop SMDR running.
\end{enumerate}
The full scale hierarchy is
$\mu_{\mathcal{N}=2} \gtrsim 37\TeV > \Lambda_{\rm CST}
> v_{\rm EW} = 246\GeV > \Lambda_F \approx 80\MeV$.


%======================================================================
\section{Open problems}
\label{sec:open}
%======================================================================

\begin{enumerate}
  \item \textbf{The K\"ahler metric.}
    Computing the K\"ahler potential on the quantum-modified moduli
    space is the master obstruction: $v_{\rm EW}$, $\tan\beta$,
    the overall scale $k$, $Z_c/Z_s$, and the scalar spectrum all
    depend on it.

  \item \textbf{The origin of $\alpha_{\rm lep}$.}
    The sector angle $\alpha_{\rm lep} = 0.2222$\,rad determines the
    lepton mass ratios but is fitted, not derived.
    All $\SU(3)$-invariant potentials have extrema at
    $\alpha = 0$ and $\pi/3$; fixing $\alpha_{\rm lep}$ requires
    non-perturbative K\"ahler effects or UV input.
    The three-instanton potential $-\cos(9\delta)$ provides one
    candidate mechanism.

  \item \textbf{The CKM integers.}
    The integers 5 and 7 in $|V_{cb}| = 1/(5R)$ and
    $|V_{ub}| = 2\sqrt{2}/(7R^3)$ are pattern identifications.
    Whether they encode Clebsch--Gordan coefficients, dimensions,
    or M-theoretic data ($84 = \binom{9}{3}$, $\pi/84$)
    is unknown.

  \item \textbf{PMNS mixing.}
    The baryonino Majorana matrix is diagonal on the mesonic branch.
    Large lepton mixing angles require non-diagonal K\"ahler
    corrections to the baryonic sector.

  \item \textbf{The quartic $\varepsilon$.}
    Computing the inter-branch coupling requires the
    non-perturbative K\"ahler metric, entangled with
    $Z_c/Z_s$ and the waterfall pattern.

  \item \textbf{The $\SO(32)$ projection.}
    Of the 496 adjoint states, 442 must be projected out.
    No mechanism has been identified.

  \item \textbf{The Casimir equation.}
    The dynamical origin of the eigenvalue equation
    $x^2 + C_2 x - C_2 = 0$ and the preferred renormalisation
    scheme remain open.
\end{enumerate}


%======================================================================
\section{Conclusions}
\label{sec:conclusion}
%======================================================================

We have presented a unified framework that begins with the
identification of Standard Model fermions as composite fermions
of a confining $\SU(3)_F$ family gauge theory---leptons as
mesinos, quarks as quarkinos---and follows the consequences
through five levels of structure:

\textbf{1.\ The composite identification} provides a single
meson matrix $M$ from which both electroweak Higgs doublets
emerge ($H_d = M^{1,2}$, $H_u = \adj(M)^{1,2}$), with no
independent Yukawa matrices.
The 9~mesinos furnish the SM lepton sector; the diquark
partners fill the symmetric $\mathbf{15}$ of $\SU(5)_f$.

\textbf{2.\ The emergent $\mathcal{N}{=}2$ doublet}
$(M_i, \adj(M)_i)$ per family, broken to $\mathcal{N}{=}1$
by a quartic coupling, generates electroweak isospin: the
direct and inverse branches produce up-type and down-type
mass hierarchies from the same charge structure.

\textbf{3.\ The charge-squared mass law} $m = \zch^2$,
uniquely selected by the requirement that $K = 2/3$ be
direction-independent in the Cartan subalgebra, produces an
energy-balance condition satisfied by the charged leptons
($0.001\%$), the meson triple $(-\pi, D_s, B)$ ($0.10\%$),
the inverse down quarks ($0.22\%$), and the quark chain
to $\sim 1\%$.
The O'Raifeartaigh model produces the seed spectrum at
$gv/m = \sqrt{3}$.

\textbf{4.\ The electroweak gauge bosons} are accommodated
through a Casimir eigenvalue equation predicting
$\sin^2\theta_W = 0.22310$ ($0.4\sigma$).

\textbf{5.\ The string/M-theory embedding}: the Diophantine
counting selects $N = 3$ generations; the composites fill
$\mathbf{24} \oplus \mathbf{15} \oplus \overline{\mathbf{15}}$
of $\SU(5)_f$, contained in the $\SO(32)$ adjoint but not
$E_8 \times E_8$.
In M-theory, mesons are M2-branes and baryons are M5-branes,
with the charge-squared mass law reflecting the membrane
dimensionality $p = 2$.

The framework has four continuous inputs for 13~observables,
with the CKM elements depending only on the self-dual ratio
$R = (5+\sqrt{21})/2$.
The most pressing open problems are the K\"ahler metric,
the origin of the sector angle $\alpha$, the CKM integers,
and the PMNS mixing pattern.


%======================================================================
\section*{Acknowledgments}
%======================================================================

During the preparation of this work the author used Claude
(Anthropic) for algebraic verification, numerical checks,
and drafting assistance.
The author reviewed and edited all content and takes full
responsibility for the publication.


%======================================================================
\begin{thebibliography}{99}

\bibitem{Chew:1962}
  G.F.~Chew,
  Rev.\ Mod.\ Phys.\ \textbf{34}, 394 (1962).

\bibitem{Miyazawa:1966}
  H.~Miyazawa,
  Prog.\ Theor.\ Phys.\ \textbf{36}, 1266 (1966).

\bibitem{CattoGursey:1985}
  S.~Catto and F.~G\"ursey,
  Nuovo Cim.\ A \textbf{86}, 201 (1985).

\bibitem{Rivero:2024epjc}
  A.~Rivero,
  Eur.\ Phys.\ J.\ C \textbf{84}, 1058 (2024),
  arXiv:2407.05397.

\bibitem{VolkovAkulov:1973}
  D.V.~Volkov and V.P.~Akulov,
  Phys.\ Lett.\ B \textbf{46}, 109 (1973).

\bibitem{FradkinShenker:1979}
  E.~Fradkin and S.H.~Shenker,
  Phys.\ Rev.\ D \textbf{19}, 3682 (1979).

\bibitem{Kaplan:1991dc}
  D.B.~Kaplan,
  Nucl.\ Phys.\ B \textbf{365}, 259 (1991).

\bibitem{Seiberg:1994bz}
  N.~Seiberg,
  Nucl.\ Phys.\ B \textbf{435}, 129 (1995),
  arXiv:hep-ph/9411149.

\bibitem{CsakiSchmaltzSkiba:1997}
  C.~Cs\'aki, M.~Schmaltz, and W.~Skiba,
  Phys.\ Rev.\ Lett.\ \textbf{78}, 799 (1997),
  arXiv:hep-th/9610139.

\bibitem{Csaki:2011xn}
  C.~Csaki, Y.~Shirman, and J.~Terning,
  Phys.\ Rev.\ D \textbf{84}, 095011 (2011),
  arXiv:1106.3074.

\bibitem{ISS:2006}
  K.~Intriligator, N.~Seiberg, and D.~Shih,
  JHEP \textbf{04}, 021 (2006),
  arXiv:hep-th/0602239.

\bibitem{Koide:1983qe}
  Y.~Koide,
  Phys.\ Lett.\ B \textbf{120}, 161 (1983).

\bibitem{Koide:1990}
  Y.~Koide,
  Mod.\ Phys.\ Lett.\ A \textbf{5}, 2319 (1990).

\bibitem{Rivero:2005}
  A.~Rivero,
  arXiv:hep-ph/0505220 (2005).

\bibitem{Rivero:2005b}
  A.~Rivero,
  arXiv:hep-ph/0510068 (2005).

\bibitem{Zenczykowski:2012}
  P.~Zenczykowski,
  Phys.\ Rev.\ D \textbf{86}, 117303 (2012).

\bibitem{Sumino:2009}
  Y.~Sumino,
  JHEP \textbf{05}, 075 (2009),
  arXiv:0812.2103.

\bibitem{HananyWitten:1997}
  A.~Hanany and E.~Witten,
  Nucl.\ Phys.\ B \textbf{492}, 152 (1997),
  arXiv:hep-th/9611230.

\bibitem{PDG:2024}
  S.~Navas \textit{et al.} (Particle Data Group),
  Phys.\ Rev.\ D \textbf{110}, 030001 (2024).

\bibitem{GattoSartoriTonin:1968}
  R.~Gatto, G.~Sartori, and M.~Tonin,
  Phys.\ Lett.\ B \textbf{28}, 128 (1968).

\bibitem{PatiSalam:1974}
  J.C.~Pati and A.~Salam,
  Phys.\ Rev.\ D \textbf{10}, 275 (1974).

\bibitem{SeibergWitten:1994}
  N.~Seiberg and E.~Witten,
  Nucl.\ Phys.\ B \textbf{426}, 19 (1994),
  arXiv:hep-th/9407087.

\bibitem{Nelson:1984}
  A.E.~Nelson,
  Phys.\ Lett.\ B \textbf{136}, 387 (1984).

\bibitem{Barr:1984}
  S.M.~Barr,
  Phys.\ Rev.\ Lett.\ \textbf{53}, 329 (1984).

\bibitem{ORaifeartaigh:1975}
  L.~O'Raifeartaigh,
  Nucl.\ Phys.\ B \textbf{96}, 331 (1975).

\bibitem{Unsal:2008}
  M.~\"Unsal,
  Phys.\ Rev.\ Lett.\ \textbf{100}, 032005 (2008),
  arXiv:0708.1772.

\bibitem{Martin:2019}
  S.P.~Martin and D.G.~Robertson,
  Phys.\ Rev.\ D \textbf{100}, 073004 (2019),
  arXiv:1907.02500.

\bibitem{GeorgiJarlskog:1979}
  H.~Georgi and C.~Jarlskog,
  Phys.\ Lett.\ B \textbf{86}, 297 (1979).

\bibitem{Townsend:1995}
  P.K.~Townsend,
  Phys.\ Lett.\ B \textbf{373}, 68 (1996),
  arXiv:hep-th/9507048.

\bibitem{BlumhoferHutter:1997}
  A.~Blumhofer and M.~Hutter,
  Nucl.\ Phys.\ B \textbf{484}, 80 (1997),
  arXiv:hep-ph/9605393.

\end{thebibliography}

\end{document}
